% ============================================================
%  PIRTM Prime Stratification Tower
%  Enforcement-Layer Architecture, Inverted Digital Twin
%  Governance Protocol, and Permissive Failure Taxonomy
%
%  PRIOR ART DEFENSIVE PUBLICATION
%  Author: Ryan Van Gelder / Multiplicity Foundation
%  First Disclosed: 2026-03-12
%  Published: 2026-04-25
% ============================================================
\documentclass[11pt,letterpaper]{article}

\usepackage[margin=1.1in]{geometry}
\usepackage{amsmath,amssymb,amsthm}
\usepackage{mathtools}
\usepackage[hidelinks]{hyperref}
\usepackage{booktabs}
\usepackage{longtable}
\usepackage{array}
\usepackage{xcolor}
\usepackage{listings}
\usepackage{caption}
\usepackage{enumitem}
\usepackage{fancyhdr}
\usepackage{microtype}
\usepackage{parskip}
\usepackage{tcolorbox}
\tcbuselibrary{skins,breakable}
\usepackage[T1]{fontenc}
\usepackage[utf8]{inputenc}
\usepackage{lmodern}

% --- Colours
\definecolor{pirtmblue}{RGB}{25,70,140}
\definecolor{pirtmgray}{RGB}{90,90,90}
\definecolor{codegreen}{RGB}{30,120,30}
\definecolor{codegray}{RGB}{130,130,130}
\definecolor{codepurple}{RGB}{120,40,140}
\definecolor{codebackground}{RGB}{248,248,248}
\definecolor{warnyellow}{RGB}{255,245,210}
\definecolor{warnborder}{RGB}{200,150,0}
\definecolor{defnblue}{RGB}{230,238,255}
\definecolor{defnborder}{RGB}{60,100,200}
\definecolor{thmgreen}{RGB}{228,255,230}
\definecolor{thmborder}{RGB}{40,160,60}

\hypersetup{
  colorlinks=true,
  linkcolor=pirtmblue,
  urlcolor=pirtmblue,
  citecolor=pirtmblue,
  pdftitle={PIRTM Prime Stratification Tower: Prior Art Defensive Publication},
  pdfauthor={Ryan Van Gelder / Citizen Gardens},
}

\pagestyle{fancy}
\fancyhf{}
\rhead{\textcolor{pirtmgray}{\small PIRTM Stratification Tower --- Defensive Publication}}
\lhead{\textcolor{pirtmgray}{\small Citizen Gardens, 2026}}
\cfoot{\thepage}
\renewcommand{\headrulewidth}{0.4pt}

% --- Code style: Python
\lstdefinestyle{python}{
  language=Python,
  backgroundcolor=\color{codebackground},
  basicstyle=\ttfamily\footnotesize,
  keywordstyle=\color{pirtmblue}\bfseries,
  stringstyle=\color{codepurple},
  commentstyle=\color{codegreen}\itshape,
  numberstyle=\tiny\color{codegray},
  breaklines=true,
  captionpos=b,
  numbers=left,
  numbersep=8pt,
  showstringspaces=false,
  frame=single,
  rulecolor=\color{pirtmgray!40},
  xleftmargin=1.8em,
  framexleftmargin=1.8em,
}

% --- Code style: MLIR
\lstdefinestyle{mlir}{
  language={},
  backgroundcolor=\color{codebackground},
  basicstyle=\ttfamily\footnotesize,
  keywordstyle=\color{pirtmblue}\bfseries,
  morekeywords={module,func,return},
  commentstyle=\color{codegreen}\itshape,
  morecomment=[l]{//},
  numberstyle=\tiny\color{codegray},
  breaklines=true,
  captionpos=b,
  numbers=left,
  numbersep=8pt,
  frame=single,
  rulecolor=\color{pirtmgray!40},
  xleftmargin=1.8em,
  framexleftmargin=1.8em,
}

% --- Code style: bash
\lstdefinestyle{bash}{
  language=bash,
  backgroundcolor=\color{codebackground},
  basicstyle=\ttfamily\footnotesize,
  keywordstyle=\color{pirtmblue}\bfseries,
  commentstyle=\color{codegreen}\itshape,
  morecomment=[l]{\#},
  stringstyle=\color{codepurple},
  breaklines=true,
  captionpos=b,
  frame=single,
  rulecolor=\color{pirtmgray!40},
  xleftmargin=1.8em,
  framexleftmargin=1.8em,
}

% --- Theorem environments
\newtheoremstyle{pirtmthm}
  {6pt}{6pt}{\itshape}{0pt}{\bfseries\color{pirtmblue}}{.}{.5em}{}
\theoremstyle{pirtmthm}
\newtheorem{definition}{Definition}[section]
\newtheorem{theorem}{Theorem}[section]
\newtheorem{proposition}{Proposition}[section]
\newtheorem{invariant}{Invariant}[section]
\theoremstyle{remark}
\newtheorem{remark}{Remark}[section]

% --- Coloured boxes
\newtcolorbox{defnbox}[1][]{
  enhanced,breakable,
  colback=defnblue,colframe=defnborder,
  fonttitle=\bfseries,title={#1},
  left=4pt,right=4pt,top=4pt,bottom=4pt,
}
\newtcolorbox{warnbox}[1][]{
  enhanced,breakable,
  colback=warnyellow,colframe=warnborder,
  fonttitle=\bfseries,title={#1},
  left=4pt,right=4pt,top=4pt,bottom=4pt,
}
\newtcolorbox{thmbox}[1][]{
  enhanced,breakable,
  colback=thmgreen,colframe=thmborder,
  fonttitle=\bfseries,title={#1},
  left=4pt,right=4pt,top=4pt,bottom=4pt,
}

% --- Title
\title{%
  \vspace{-1.5em}
  {\Large\textbf{\textcolor{pirtmblue}{PIRTM Prime Stratification Tower}}}\\[4pt]
  {\large Enforcement-Layer Architecture, Inverted Digital Twin\\
   Governance Protocol, and Permissive Failure Taxonomy}\\[6pt]
  {\normalsize\textbf{Prior Art Defensive Publication}}
}
\author{Citizen Gardens \\ \texttt{The foundation of Multiplicity}}
\date{April 25, 2026}

\begin{document}
\maketitle
\thispagestyle{fancy}

\begin{abstract}
This document constitutes a \emph{prior art defensive publication} disclosing the
\textbf{PIRTM Prime Stratification Tower}: a formally specified, enforcement-layered
software architecture for the Prime Interval Recursion Theory Machine (PIRTM).
The disclosed system combines four independently novel contributions:
(1)~a prime-indexed type stratification system grounded in $p$-adic valuation theory,
in which each operational stratum is labeled by a prime $p_n$ and cross-stratum
call validity is governed by a prime-index successor predicate;
(2)~an \emph{enforcement layer separation} principle assigning each invariant
exclusively to the surface architecturally capable of checking it;
(3)~an \emph{Inverted Digital Twin} (IDT) governance mechanism in which a
ground-stratum oracle at $\bot$ actively verifies the stratum-$p_1$ verifier
pass on every commit;
and (4)~a \emph{permissive failure taxonomy} identifying five structurally distinct
modes by which an enforcement surface produces false assurance without any
diagnostic.
The publication establishes a publicly dated record as of 2026-03-12,
placing all described techniques in the public domain.
\end{abstract}
\newpage
\tableofcontents
\newpage

%%------------------------------------------------------------
\section{Executive Summary}
%%------------------------------------------------------------

The PIRTM Prime Stratification Tower is a formal governance architecture for
compiler-integrated DSL development within MLIR.
Its central problem is \emph{permissive failure}: a condition in which an
enforcement surface nominally checks an invariant but silently accepts violations,
producing false assurance indistinguishable from genuine compliance.

The disclosed architecture addresses this through four interlocking mechanisms:

\begin{enumerate}[leftmargin=1.8em]
  \item \textbf{Prime-indexed stratification.}
    Each PIRTM operation is assigned a stratum from the sequence
    $\{p_1, p_2, p_3, \ldots\} = \{2, 3, 5, 7, 11, \ldots\}$ via
    \texttt{pirtm.stratum}.
    Cross-stratum call validity uses the prime-index successor predicate
    $\mathrm{idx}(s(A)) \leq \mathrm{idx}(s(B)) + 1$, where $A$ is the
    consumer and $B$ is the producer.

  \item \textbf{Enforcement layer separation.}
    Each invariant is assigned to exactly one enforcement layer: MLIR verifier
    passes for stratum annotation and call direction; CI shell scripts for
    path scope and ratio; issue trackers for proof obligations.

  \item \textbf{Inverted Digital Twin (IDT) governance.}
    A committed \texttt{.mlir} fixture with a paired \texttt{.expected} oracle
    file causes CI to fail when the verifier pass deviates from the oracle in
    either direction.

  \item \textbf{Permissive failure taxonomy.}
    Five structurally distinct failure forms are catalogued, each with a
    detection mechanism and a committed-artifact remediation protocol.
\end{enumerate}

The architecture was developed through thirteen iterative critique-and-correction
rounds on 2026-03-12, with each round surfacing a new permissive failure instance.
The system converges when every invariant is enforced by exactly one mechanism,
every mechanism is tested by an independent oracle, and every named-but-not-committed
fix is tracked in a dated, CI-visible artifact.

%%------------------------------------------------------------
\section{Background and Motivation}
%%------------------------------------------------------------

\subsection{The PIRTM System}

The Prime Interval Recursion Theory Machine (PIRTM) is a formal computation
model developed under Multiplicity Theory.
Mathematical objects are reinterpreted as recursively generated patterns of
prime-labeled interactions; each level of the computation tower corresponds to
a prime $p_n$; behavior is preserved through recursive feedback loops across scales.
PIRTM is implemented as an MLIR dialect with operations carrying attributed
prime-indexed strata.

\subsection{The Decorative Governance Problem}

A system exhibits \emph{decorative governance} when its policy documents,
naming conventions, and structural rules appear to enforce architectural
constraints but do not.
The PIRTM development process identified this as the root of a repeated pattern:
corrective code itself instantiated the failure mode it was designed to prevent,
because the corrective layer was written without an active machine check below it.

\subsection{MLIR Context}

MLIR's Python bindings allow verifier passes to be prototyped without
C++/TableGen dialect registration, using
\texttt{Context(allow\_unregistered\_dialects=True)} to inspect named attributes.
The key API is \texttt{IntegerAttr(attr).value}, which extracts \texttt{int64}
directly from the MLIR C API.
This avoids the \texttt{int(str(attr))} permissive failure, which silently raises
\texttt{ValueError} on every correctly-annotated op because integer attribute
serialization includes the type suffix (e.g., \texttt{"5 : i64"}, not \texttt{"5"}).

%%------------------------------------------------------------
\section{Mathematical Foundations}
%%------------------------------------------------------------

\subsection{Prime-Indexed Stratum Domain}

\begin{definition}[Prime Index Map]
Let $\mathcal{P} = \{p_1, p_2, p_3, \ldots\} = \{2, 3, 5, 7, 11, \ldots\}$
be the ordered prime sequence.
The \emph{prime index map} is
\[
  \mathrm{idx} : \mathcal{P} \to \mathbb{N}, \qquad \mathrm{idx}(p_n) = n.
\]
The \emph{prime successor} is
$\mathrm{succ}_{\mathcal{P}}(p_n) = p_{n+1}$.
Implementation: \texttt{PRIME\_IDX = \{2:1, 3:2, 5:3, 7:4, 11:5, \ldots\}}.
\end{definition}

\begin{definition}[Extended Stratum Domain]
$\mathcal{S} = \{\bot\} \cup \mathcal{P}$,
where $\bot$ denotes the exterior ground element for Python-bound operations
outside the prime ring.
\end{definition}

\subsection{\texorpdfstring{$p$}{p}-Adic Grounding}

For a fixed prime $p$, the $p$-adic valuation $v_p : \mathbb{Q}_p^* \to \mathbb{Z}$
satisfies $v_p(p^n) = n$ for any $n \in \mathbb{Z}$, including $n < 0$.
Elements in $\mathbb{Q}_p \setminus \mathbb{Z}_p$ have $v_p < 0$.
The $\bot$ ground corresponds to this negative-valuation regime.
PIRTM operations at level $n$ live in $\mathbb{Z}_p$ with $v_p \geq 0$.

\begin{warnbox}[Notational Convention on $-1$]
The label $-1$ for $\bot$ used in early development is operational shorthand
for a build-dependency rank, not a literal $p$-adic valuation index.
The formal system uses $\bot$ throughout.
\end{warnbox}

\subsection{The ValidCall Predicate}

\begin{definition}[ValidCall]
For operations $A$ (consumer) and $B$ (producer), the call is valid iff:
\[
  \mathrm{ValidCall}(A, B) \equiv
  \begin{cases}
    \top
      & \text{if } s(A) = \bot \text{ and } B \text{ is a compiled PIRTM artifact}\\
    \bot
      & \text{if } s(B) = \bot \text{ and } A \in \mathcal{O}\\
    \mathrm{idx}(s(A)) \leq \mathrm{idx}(s(B)) + 1
      & \text{if } s(A), s(B) \in \mathcal{P}.
  \end{cases}
\]
\textit{A consumer's prime index cannot exceed its producer's prime index by more than 1.}
\end{definition}

\begin{warnbox}[Index vs.\ Value Arithmetic]
The $+1$ operates on prime \emph{indices} ($1, 2, 3, \ldots$), not prime
\emph{values} ($2, 3, 5, 7, \ldots$).
Example: $s(B)=p_3=5$ gives $\mathrm{idx}(5)+1=4$, so consumer is permitted
up to $p_4=7$ and rejected at $p_5=11$ because $\mathrm{idx}(11)=5>4$.
The expression $s(B)+1=6$ is not a prime and is semantically wrong.
All implementations use \texttt{PRIME\_IDX[s(B)]+1}, never \texttt{s(B)+1}.
\end{warnbox}

\subsection{Scoped Ratio Metric}

Let $L = |\text{.mlir lines in \texttt{src/logic/}}|$ and
$P = |\text{.py lines in \texttt{src/logic/}}|$ (generated files excluded).
\[
  r_k =
  \begin{cases}
    \text{undefined (exit 2)} & \text{if } L + P = 0\\[4pt]
    \dfrac{L}{L + P} & \text{otherwise}
  \end{cases}
  \qquad \text{CI fails iff } r_k < \tau_k,
\]
where $\tau_1 = 0.50$, $\tau_2 = 0.65$, $\tau_3 = 0.80$
are committed constants in \texttt{measure\_ratio.sh}.

\subsection{IDT Oracle Predicate}

Let $\mathcal{D} : \mathcal{O} \to 2^{\Sigma^*}$ map each op to its diagnostic
message strings.
The oracle matching predicate is:
\begin{align*}
  \mathrm{OracleMatch}(\hat{\mathcal{D}},\Omega) \equiv &
    \Bigl(\forall e\in\Omega,\;
    \exists\,act\in\textstyle\bigcup_{op}\hat{\mathcal{D}}(op):
    e \subseteq act\Bigr)\\
  &\wedge
    \Bigl(\forall\,act\in\textstyle\bigcup_{op}\hat{\mathcal{D}}(op),\;
    \texttt{": error:"}\in act \implies \exists\,e\in\Omega: e\subseteq act\Bigr),
\end{align*}
where $\subseteq$ denotes substring containment.
This is path-independent (weaker than equality) and op-name-anchored
(stronger than count matching).

\subsection{Enforcement Tower Stability}

\[
  \forall k \geq 0,\quad
  \|\phi_{\mathrm{enforcer}}^k(E) - E^*\|_{p_1} \to 0,
\]
where $E^*$ is the oracle error set and $\phi_{\mathrm{enforcer}}$ is the
pass applied iteratively.
The $\bot$ oracle is stable by definition; no meta-meta-verifier is needed.

\subsection{Session Debt}

\begin{definition}[Form-4 Debt and Escalation]
\[
  \mathrm{DEBT}(\mathcal{S}_t) =
  \{f \mid f \text{ named in session } \mathcal{S}_t,\;
    \nexists\,\text{committed artifact for } f \text{ in } \mathcal{S}_t\}.
\]
\[
  \mathrm{Escalates}(f) = \mathrm{Named}(f) + 21\,\text{days}.
\]
Session closure requires $\mathrm{DEBT}(\mathcal{S}_t)
\subseteq \texttt{SESSION\_DEBT.md}$ at close of $\mathcal{S}_t$.
\end{definition}

%%------------------------------------------------------------
\section{Permissive Failure Taxonomy}
%%------------------------------------------------------------

A \emph{permissive failure} is a condition in which an enforcement surface
accepts a violation silently, producing false assurance.

\begin{center}
\renewcommand{\arraystretch}{1.35}
\begin{tabular}{@{}>{\bfseries}p{2.6cm}p{4.8cm}p{4.6cm}@{}}
  \toprule
  Form & Instance & False Assurance \\
  \midrule
  Wrong layer
    & \texttt{POLICY.md} enforcing a CI-only invariant
    & Document looks like enforcement \\
  Wrong mechanism
    & \texttt{int(str(attr))}, \texttt{git grep},
      \texttt{mlirAttributeIsAInteger} from Python,
      \texttt{assert} under \texttt{-O}
    & Code looks like it runs correctly \\
  Wrong direction
    & Oracle predicate inverted; ValidCall inequality inverted
    & Tests look like they check the right thing \\
  Named only
    & Substring matching described in prose, equality committed
    & Fix looks present in the record \\
  Unvalidated spec
    & \texttt{open(sys.argv)} opens a list object;
      \texttt{/dev/null} as MLIR smoke input;
      trailing periods in ISO-8601 \texttt{Escalates:} lines
    & Proposed code looks runnable \\
  \bottomrule
\end{tabular}
\end{center}

\begin{thmbox}[Meta-Pattern]
All five forms share a single root: an enforcement surface was written without
an active machine check at the level below it.
The IDT oracle at $\bot$ detects Forms~2, 3, and~4.
Form~1 is detected by the enforcement layer table.
Form~5 is detected by \texttt{check\_stub.py} and input-domain validation.
The taxonomy itself lacks a machine oracle; \texttt{SESSION\_DEBT.md} is the
manual substitute.
\end{thmbox}

%%------------------------------------------------------------
\section{Enforcement Layer Architecture}
%%------------------------------------------------------------

\subsection{Separation Principle}

Every invariant is assigned to exactly one enforcement layer.
\texttt{mlir-opt} is file-path-agnostic; it cannot enforce path-scope invariants.
CI shell scripts have no access to the IR use-def graph.
Issue trackers have neither.
Assigning an invariant to the wrong layer is Form~1 permissive failure.

\begin{center}
\renewcommand{\arraystretch}{1.3}
\begin{tabular}{@{}>{\small\ttfamily}p{4.4cm}p{3.0cm}p{3.6cm}@{}}
  \toprule
  \textrm{Invariant} & \textrm{Layer} & \textrm{Artifact} \\
  \midrule
  Stratum annotation present         & MLIR verifier & \texttt{pirtm\_verifier.py} \\
  PRIME\_IDX[s(A)] <= PRIME\_IDX[s(B)]+1 & MLIR verifier & \texttt{pirtm\_verifier.py} \\
  pirtm.extern path-scoped           & CI script & \texttt{check\_extern\_scope.sh} \\
  Ratio $\geq \tau_k$                & CI script & \texttt{measure\_ratio.sh} \\
  Generated markers present          & CI script & \texttt{check\_generated\_markers.sh} \\
  Sufficiency tracked                & Issue tracker & Dated issue, canonical repo \\
  \bottomrule
\end{tabular}
\end{center}

\subsection{Full Stratum Tower}

\begin{center}
\renewcommand{\arraystretch}{1.3}
\begin{tabular}{@{}>{\bfseries}p{1.3cm}p{3.8cm}p{2.2cm}p{3.5cm}@{}}
  \toprule
  Stratum & Artifact & Governs & Failure if Absent \\
  \midrule
  $\bot$   & Fixture + oracle              & Verifier at $p_1$    & Correctness unverifiable \\
  $p_0$    & CI scripts                    & Ratio \& paths       & Invariants decorative \\
  $p_1$    & \texttt{pirtm\_verifier.py}   & PIRTM IR at $p_2$    & Violations accepted \\
  $p_1^*$  & Dated sufficiency issue       & Proof obligations    & Stability ungrounded \\
  $p_2$    & Governance twin (PIRTM IR)    & DSL surface          & No reference impl. \\
  $p_3$    & DSL surface                   & Authoring exp.       & Python dialect persists \\
  \bottomrule
\end{tabular}
\end{center}

%%------------------------------------------------------------
\section{The Inverted Digital Twin (IDT)}
%%------------------------------------------------------------

The \emph{Inverted Digital Twin} places the oracle at $\bot$ so that the
higher-stratum pass ($p_1$) cannot pass CI without satisfying it.
The IDT oracle consists of:
\begin{enumerate}[leftmargin=1.8em]
  \item A committed \texttt{.mlir} fixture with four diagnostic-bearing ops;
  \item A paired \texttt{.expected} file with op-name-anchored error substrings;
  \item A \texttt{check\_oracle.py} CI script that fails on deviation in either direction.
\end{enumerate}

\subsection{Four-Op Fixture Specification}

\begin{center}
\renewcommand{\arraystretch}{1.3}
\begin{tabular}{@{}>{\ttfamily\small}p{3.6cm}p{1.8cm}p{5.8cm}@{}}
  \toprule
  \textrm{Op name} & Expected & Purpose \\
  \midrule
  pirtm.baseline       & 0 errors & Walk ran; annotation check passes standalone \\
  pirtm.unannotated    & 1 error  & Missing \texttt{pirtm.stratum} detected \\
  pirtm.good\_consumer & 0 errors & $p_4=7$ from $p_3=5$: $\mathrm{idx}(7)=4\leq3+1$ \\
  pirtm.bad\_consumer  & 1 error  & $p_5=11$ from $p_3=5$: $\mathrm{idx}(11)=5>3+1$ \\
  \bottomrule
\end{tabular}
\end{center}

Four-state diagnostic signature:
\begin{itemize}[leftmargin=1.8em]
  \item 0 errors: walk never ran (permissive)
  \item 1 error on unannotated only: ValidCall walk broken (permissive)
  \item 1 error on bad\_consumer only: annotation check broken (permissive)
  \item 2 errors on correct ops: pass correct
  \item 3+ errors: over-rejection; idx arithmetic wrong
\end{itemize}

%%------------------------------------------------------------
\section{Code Listings}
%%------------------------------------------------------------

\subsection{pirtm\_verifier.py --- Complete Implementation Skeleton}

\begin{lstlisting}[style=python,
  caption={pirtm\_verifier.py --- Verifier with get\_stratum, check\_valid\_call, walk, and entrypoint}]
# pirtm_verifier.py
# Invariant: PRIME_IDX[s(consumer)] <= PRIME_IDX[s(producer)] + 1
# where s(op) = IntegerAttr(op.attributes["pirtm.stratum"]).value
# PRIME_IDX = {2:1, 3:2, 5:3, 7:4, 11:5, 13:6, 17:7, 19:8, 23:9, 29:10}
# Convention: every op.emit_error includes op.name as first token.

from mlir.ir import Context, Module, IntegerAttr, MLIRError
from pathlib import Path
import sys

PRIMES = [2, 3, 5, 7, 11, 13, 17, 19, 23, 29]
PRIME_IDX = {p: i+1 for i, p in enumerate(PRIMES)}

def get_stratum(op):
    # Returns prime value in PRIME_IDX, or None after emitting op.emit_error.
    attr = op.attributes.get("pirtm.stratum")
    if attr is None:
        op.emit_error(f"{op.name}: missing pirtm.stratum annotation")
        return None
    if not isinstance(attr, IntegerAttr):
        op.emit_error(
            f"{op.name}: pirtm.stratum must be IntegerAttr; "
            f"got {type(attr).__name__}"
        )
        return None
    val = IntegerAttr(attr).value   # direct i64 extraction -- no str()
    if val not in PRIME_IDX:
        op.emit_error(
            f"{op.name}: stratum value {val} is not a valid prime; "
            f"valid: {PRIMES}"
        )
        return None
    return val
    # Exhaustiveness proof: non-None iff (1) attr present, (2) IntegerAttr,
    # (3) val in PRIME_IDX.  Downstream assert is dead code; guard is here.

def check_valid_call(consumer_op, producer_op):
    # None = get_stratum already emitted error; skip ValidCall, continue walk.
    # "skip" != "accept": the diagnostic is already in the error stream.
    cs = get_stratum(consumer_op)
    ps = get_stratum(producer_op)
    if cs is None or ps is None:
        return
    if PRIME_IDX[cs] > PRIME_IDX[ps] + 1:
        consumer_op.emit_error(
            f"{consumer_op.name}: stratum {cs} "
            f"(idx={PRIME_IDX[cs]}) cannot consume from "
            f"stratum {ps} (idx={PRIME_IDX[ps]}); "
            f"max allowed consumer idx: {PRIME_IDX[ps]+1}"
        )

def walk(module):
    for op in module.body.operations:
        for region in op.regions:
            for block in region.blocks:
                for inner_op in block.operations:
                    get_stratum(inner_op)   # annotation check
                    for operand in inner_op.operands:
                        owner = operand.owner
                        if hasattr(owner, 'operation'):
                            check_valid_call(inner_op, owner.operation)
                        # BlockArgument: policy per AGENTS.md

if __name__ == "__main__":
    if len(sys.argv) < 2:
        print("Usage: pirtm_verifier.py <fixture.mlir>")
        sys.exit(1)
    fixture = Path(sys.argv[1]).resolve()
    with Context() as ctx:
        ctx.allow_unregistered_dialects = True
        try:
            with open(fixture) as f:
                module = Module.parse(f.read())
        except FileNotFoundError:
            print(f"file not found: {fixture}", file=sys.stderr)
            sys.exit(2)
        except MLIRError as e:
            print(f"parse failure: {e}", file=sys.stderr)
            sys.exit(2)    # exit 2: parse failure, not verifier failure
        walk(module)
\end{lstlisting}

\subsection{check\_oracle.py --- IDT Oracle Runner}

\begin{lstlisting}[style=python,
  caption={check\_oracle.py --- Full input-domain validation; substring oracle matching}]
# check_oracle.py -- IDT comparison mechanism
# Three guards validate input domain before any logic executes.
from pathlib import Path
import subprocess, sys

VERIFIER = Path(__file__).parent / "pirtm_verifier.py"

# Guard 1: argument count
if len(sys.argv) < 2:
    print("Usage: check_oracle.py <fixture.mlir>")
    sys.exit(1)

# Guard 2: extension on raw string
if not sys.argv[1].endswith(".mlir"):
    print(f"Input must be a .mlir file; got {sys.argv[1]}")
    sys.exit(1)

# Resolve once; all downstream paths absolute
fixture     = Path(sys.argv[1]).resolve()
oracle_file = fixture.with_suffix(".expected")

def load_lines(path):
    try:
        with open(path) as f:
            return [l.strip() for l in f if l.strip()]
    except FileNotFoundError:
        print(f"ORACLE FILE MISSING: {path}")
        sys.exit(1)

# Guard 3: oracle existence BEFORE subprocess
oracle = load_lines(oracle_file)

result = subprocess.run(
    [sys.executable, str(VERIFIER), str(fixture)],
    capture_output=True, text=True
)
actual = result.stderr.splitlines()

# Substring matching: path-independent, op-name-anchored.
# ": error:" colon-space prefix excludes "error:" in note body text.
missing    = [e for e in oracle if not any(e in a for a in actual)]
unexpected = [a for a in actual if ": error:" in a
              and not any(e in a for e in oracle)]

if missing or unexpected:
    [print(f"MISSING:    '{e}'") for e in missing]
    [print(f"UNEXPECTED: '{a}'") for a in unexpected]
    sys.exit(1)
print("IDT PASS")
\end{lstlisting}

\subsection{check\_stub.py}

\begin{lstlisting}[style=python, caption={check\_stub.py --- Smoke test for verifier stub}]
# check_stub.py -- smoke test: stub parses stub_smoke.mlir cleanly.
# Does NOT require zero stderr; MLIR context init may emit informational output.
from pathlib import Path
import subprocess, sys

VERIFIER = Path(__file__).parent / "pirtm_verifier.py"
SMOKE    = Path(__file__).parent / "stub_smoke.mlir"

result = subprocess.run(
    [sys.executable, str(VERIFIER), str(SMOKE)],
    capture_output=True, text=True
)
if result.returncode != 0:
    print(f"STUB FAIL: exit {result.returncode}\n{result.stderr}")
    sys.exit(1)
if ": error:" in result.stderr:    # mirrors check_oracle.py filter
    print(f"STUB FAIL: error in stderr:\n{result.stderr}")
    sys.exit(1)
print("STUB PASS")
\end{lstlisting}

\subsection{CI Shell Scripts}

\begin{lstlisting}[style=bash,
  caption={measure\_ratio.sh --- Sprint ratio metric with bootstrap guard}]
#!/bin/bash
# Bootstrap guard: exit 2 if empty denominator (not a policy failure).
# Generated files excluded by name pattern and # GENERATED marker.
SPRINT=${SPRINT:-1}
case $SPRINT in
  1) TAU=0.50 ;; 2) TAU=0.65 ;; 3) TAU=0.80 ;;
  *) echo "Unknown sprint $SPRINT"; exit 1 ;;
esac
PY=$(find src/logic -name '*.py' \
  ! -name '*_generated.py' ! -name '*_pb2.py' ! -name '*.pyi' \
  | xargs grep -rL "# GENERATED" 2>/dev/null | wc -l)
ML=$(find src/logic -name '*.mlir' | wc -l)
TOTAL=$((PY + ML))
[ "$TOTAL" -eq 0 ] && echo "ratio=undefined (bootstrap)" && exit 2
RATIO=$(echo "scale=4; $ML / $TOTAL" | bc)
echo "ratio=$RATIO (threshold=$TAU, sprint=$SPRINT)"
(( $(echo "$RATIO < $TAU" | bc -l) )) && exit 1
exit 0
\end{lstlisting}

\begin{lstlisting}[style=bash,
  caption={check\_extern\_scope.sh --- Path-scope enforcement for pirtm.extern}]
#!/bin/bash
# grep -r for CI/local consistency; git grep misses untracked files
if grep -r 'pirtm.extern' src/logic/ 2>/dev/null | grep -q .; then
  echo "VIOLATION: pirtm.extern found in src/logic/; must be src/runner/ only"
  exit 1
fi
exit 0
\end{lstlisting}

\begin{lstlisting}[style=bash,
  caption={check\_generated\_markers.sh --- Generator discipline enforcement}]
#!/bin/bash
# Bootstrap guard: exit 0 if no generated files (mirrors measure_ratio.sh).
files=$(find src/logic -name '*_generated.py' -o -name '*_pb2.py' 2>/dev/null)
[ -z "$files" ] && exit 0
echo "$files" | xargs grep -rL "# GENERATED" | grep . && exit 1
exit 0
\end{lstlisting}

\subsection{MLIR Fixture and Oracle Files}

\begin{lstlisting}[style=mlir,
  caption={test\_stratum\_violation.mlir --- Four-op IDT oracle fixture}]
// test_stratum_violation.mlir
// Oracle: exactly 2 errors (pirtm.unannotated, pirtm.bad_consumer)
module {
  // Op 1: baseline -- confirms walk ran; annotation passes standalone
  %0 = "pirtm.baseline"() {pirtm.stratum = 5 : i64} : () -> i64

  // Op 2: EXPECT error -- missing pirtm.stratum annotation
  %1 = "pirtm.unannotated"(%0) : (i64) -> i64

  // Op 3: EXPECT clean -- p4=7 from p3=5: idx(7)=4 <= idx(5)+1=4
  %2 = "pirtm.good_consumer"(%0) {pirtm.stratum = 7 : i64} : (i64) -> i64

  // Op 4: EXPECT error -- p5=11 from p3=5: idx(11)=5 > idx(5)+1=4
  %3 = "pirtm.bad_consumer"(%0) {pirtm.stratum = 11 : i64} : (i64) -> i64
}
\end{lstlisting}

\noindent
\texttt{test\_stratum\_violation.expected} (written after message format confirmed on Day~2):
\begin{verbatim}
pirtm.unannotated: missing pirtm.stratum annotation
pirtm.bad_consumer: stratum 11
\end{verbatim}

\noindent
\texttt{stub\_smoke.mlir} (minimal valid MLIR for smoke test):
\begin{verbatim}
module {}
\end{verbatim}

%%------------------------------------------------------------
\section{Dependency Stack and Precondition Chain}
%%------------------------------------------------------------

\begin{center}
\renewcommand{\arraystretch}{1.35}
\begin{tabular}{@{}>{\bfseries}p{1.3cm}p{3.4cm}p{2.8cm}p{3.5cm}@{}}
  \toprule
  Stratum & Artifact & Precondition & Acceptance Criterion \\
  \midrule
  $p_{-1}$
    & Dirs + READMEs
    & None
    & \texttt{find src/logic} parseable \\
  $p_0$
    & 3 CI scripts + migrated file
    & Dirs exist; denom nonzero
    & CI exits per exit-code convention \\
  $p_1$
    & Verifier + IDT oracle
    & \texttt{mlir-opt} installed
    & IDT PASS on fixture \\
  $p_1^*$
    & Dated sufficiency issue
    & Canonical repo declared
    & Issue URL at correct repo \\
  $p_2$
    & Governance twin (PIRTM IR)
    & Pass at $p_1$ accepts twin
    & Twin accepted by pass \\
  $p_3$
    & DSL surface
    & Twin as reference impl.
    & DSL round-trips twin IR \\
  \bottomrule
\end{tabular}
\end{center}

%%------------------------------------------------------------
\section{SESSION\_DEBT.md Protocol}
%%------------------------------------------------------------

\texttt{SESSION\_DEBT.md} is the operational substitute for a meta-oracle on
specification prose, converting Form-4 failures from conversation history into
committed, CI-visible artifacts with dated escalation clocks.

Protocol rules:
\begin{enumerate}[leftmargin=1.8em]
  \item Log every Form-4 before session close with \texttt{Named: YYYY-MM-DD}.
  \item Remove an entry only when its artifact is referenced in a commit message.
  \item $\mathrm{Escalates}(f) = \mathrm{Named}(f) + 21$ days.
  \item Escalated items are blocking dependencies, not technical debt.
\end{enumerate}

Machine parseability: \texttt{Escalates:} lines use ISO-8601 with no trailing
punctuation, enabling extraction via
\texttt{grep -oE "[0-9]\{4\}-[0-9]\{2\}-[0-9]\{2\}"}.

\noindent Initial open items (named 2026-03-12, escalate 2026-04-02):
\begin{itemize}[leftmargin=1.8em,itemsep=2pt]
  \item Custom DiagnosticHandler to replace \texttt{": error:"} filter
  \item BlockArgument stratum policy ($\bot$ vs.\ enclosing op stratum)
  \item Sufficiency proof obligation for $\mathrm{SVC}_{p_n}$
  \item Meta-oracle for specification prose
  \item \texttt{SESSION\_DEBT.md} CI enforcement binding (\texttt{check\_debt.py})
\end{itemize}

%%------------------------------------------------------------
\section{Claims for Prior Art}
%%------------------------------------------------------------

This publication establishes prior art as of 2026-03-12.
No patent rights are claimed.
All described techniques are placed in the public domain.

\begin{enumerate}[leftmargin=1.8em,label=\textbf{Claim \arabic*.},itemsep=4pt]
  \item A method of assigning software operations to stratification levels using
    the prime sequence, wherein cross-stratum call validity is governed by
    $\mathrm{idx}(s(A)) \leq \mathrm{idx}(s(B)) + 1$ on prime indices, not values.

  \item A compiler verifier pass using MLIR Python bindings with
    \texttt{allow\_unregistered\_dialects=True}, wherein every error path
    routes to \texttt{op.emit\_error} with the operation name as first token,
    and stratum values are extracted via \texttt{IntegerAttr(attr).value}.

  \item An Inverted Digital Twin governance mechanism comprising a committed
    \texttt{.mlir} fixture, a paired \texttt{.expected} oracle file, and a CI
    comparison script that fails on deviation in either direction.

  \item An enforcement layer separation architecture assigning each invariant
    exclusively to one of: MLIR verifier pass, CI shell script, or issue tracker.

  \item A five-form permissive failure taxonomy (wrong layer, wrong mechanism,
    wrong direction, named-only, unvalidated spec) with a committed-artifact
    remediation protocol for each form.

  \item A session debt tracking protocol with ISO-8601 escalation dates,
    a three-calendar-week clock, and a CI-enforced deadline mechanism.

  \item A scoped ratio metric $r_k = L/(L+P)$ with sprint-locked thresholds,
    a bootstrap guard (exit~2 for empty denominator), and generated-file exclusion.

  \item A ground-stratum oracle at $\bot$ as the immutable fixed point of the
    enforcement tower, such that recursive stability requires no meta-meta-verifier.

  \item A path-resolved IDT comparison script with three-guard input-domain
    validation (argument count, extension, oracle existence) producing structured
    diagnostics for all boundary conditions.

  \item A \texttt{PRIME\_IDX} precomputed dictionary with the construction
    invariant that any non-\texttt{None} return from \texttt{get\_stratum}
    is a key in \texttt{PRIME\_IDX}, eliminating redundant downstream assertions.
\end{enumerate}

%%------------------------------------------------------------
\section{Day-by-Day Execution Timeline}
%%------------------------------------------------------------

\begin{center}
\renewcommand{\arraystretch}{1.35}
\begin{tabular}{@{}>{\bfseries}p{0.9cm}p{1.6cm}p{8.5cm}@{}}
  \toprule
  Day & Est. & Tasks and Acceptance Criteria \\
  \midrule
  0 & 3--4 hrs
    & Create dirs; migrate one logic file; write \texttt{CANONICAL\_REPO};
      commit seven-file stub manifest.
      Confirm: STUB PASS + ORACLE FILE MISSING locally before push. \\
  1 & 2--3 hrs
    & Commit three CI scripts together.
      Run \texttt{measure\_ratio.sh}; confirm exit~0 or~2 (never~1 at this stage). \\
  2 & 4--6 hrs
    & Write \texttt{.expected} after confirming message format.
      Implement walk. Confirm IDT PASS (exactly 2 errors, correct ops). \\
  3 & 1--2 hrs
    & Commit \texttt{pirtm\_authoring\_policy.mlir} stub.
      Commit BlockArgument policy to \texttt{AGENTS.md}. \\
  4 & 1 hr
    & Extend verifier to full use-def walk; confirm rejection of known violation. \\
  5 & 30 min
    & File sufficiency issue at \texttt{CANONICAL\_REPO} URL.
      Issue URL is now a verifiable artifact. \\
  6--10 & ongoing
    & Governance twin rewrite. Blocked until Day~2 pass is green.
      C++ dialect stub at Day~7+. \\
  \bottomrule
\end{tabular}
\end{center}

%%------------------------------------------------------------
\section{Open Issues and Future Work}
%%------------------------------------------------------------

\begin{enumerate}[leftmargin=1.8em,itemsep=4pt]
  \item \textbf{Custom DiagnosticHandler.}
    Replace the \texttt{": error:"} substring filter in \texttt{check\_oracle.py}
    with a registered handler capturing diagnostics as structured objects.

  \item \textbf{BlockArgument stratum policy.}
    Decide before finalizing the fixture: assign $\bot$ or inherit enclosing
    op stratum.
    The oracle $E(\mathcal{F})$ changes with this choice.

  \item \textbf{Sufficiency proof for $\mathrm{SVC}_{p_n}$.}
    A Lean or Coq proof that linear $\phi$ over $\mathbb{Z}_{p_n}$ is a
    contraction under the $p$-adic ultrametric.
    Recursive stability claims in the twin rewrite are blocked until this closes.

  \item \textbf{C++ dialect stub.}
    Register \texttt{pirtm.stratum} and \texttt{pirtm.extern} in a compiled
    dialect, enabling lit \texttt{//~expected-error~\{\{\ldots\}\}} syntax.

  \item \textbf{Extension guard ordering.}
    Refactor \texttt{check\_oracle.py}: resolve $\to$ \texttt{.suffix~==~".mlir"},
    handling symlinks with non-\texttt{.mlir} names.

  \item \textbf{Meta-oracle for specification prose.}
    A formal specification language or automated review protocol applying the
    IDT principle to specification documents.
\end{enumerate}

%%------------------------------------------------------------
\section{Publication Statement}
%%------------------------------------------------------------

This document is published as a defensive disclosure to establish prior art
under 35~U.S.C.~102(a)(1) and equivalent statutes in applicable jurisdictions.
The disclosed subject matter is placed in the public domain as of 2026-03-12.
No patent rights are claimed by or reserved to the authors.
All described algorithms, architectures, and protocols are freely available
for use, implementation, and further development without restriction.

The Multiplicity Foundation asserts authorship and first-disclosure priority
for all claims in Section~10 as of 2026-03-12.
This publication may be cited as prior art in any patent examination or
validity proceeding involving the described subject matter.

\vspace{2em}
\noindent\rule{\textwidth}{0.4pt}
\vspace{0.5em}

\noindent
\textbf{Multiplicity Foundation / Citizen Gardens 501(c)(3)} \\
Bridgeville, PA, USA \\[2pt]
\textbf{First Disclosed:} 2026-03-12 \qquad
\textbf{Published:} 2026-04-25 \\
\textbf{Author:} Ryan Van Gelder

\appendix}
\section{Operator Norm Bounds on the PIRTM Step Map}
\label{sec:opnorm}
% ════════════════════════════════════════════════════════════

\subsection{Setup and Notation}

Recall the PIRTM recurrence from the main text:
\begin{equation}
  X_{t+1} = P\!\left(\Xi\, X_t + \Lambda\, T(X_t) + G_t\right),
  \label{eq:A-recurrence}
\end{equation}
where $X_t \in \R^n$, $\Xi, \Lambda \in \R^{n \times n}$,
$T(x) = \sigma(x) = (1+e^{-x})^{-1}$ component-wise,
$G_t \in \R^n$, and $P = \operatorname{clip}(\cdot,-1,1)$ component-wise.

All vector norms are the Euclidean norm $\norm{\cdot}_2$ unless stated
otherwise. All matrix norms are the induced spectral norm
$\norm{A}_2 = \sigma_{\max}(A)$ (largest singular value).
We write $r(A) = \max\{|\lambda| : \lambda \in \spec(A)\}$ for the
spectral radius.

\subsection{Lipschitz Bounds on the Sub-operators}

\begin{lemma}[Sigmoid Lipschitz constant]
\label{lem:sigmoid-lip}
The sigmoid $\sigma : \R \to (0,1)$, $\sigma(x) = (1+e^{-x})^{-1}$,
satisfies $\sigma'(x) = \sigma(x)(1-\sigma(x)) \leq \tfrac{1}{4}$
for all $x \in \R$, with equality at $x = 0$.
Consequently, the component-wise map $T : \R^n \to \R^n$
satisfies:
\[
  \norm{T(x) - T(y)}_2 \leq \tfrac{1}{4}\norm{x-y}_2
  \quad \forall\, x, y \in \R^n.
\]
Hence $\Lip(T) = \tfrac{1}{4}$.
\end{lemma}

\begin{proof}
By the Mean Value Theorem, for each component $i$:
\[
  |\sigma(x_i) - \sigma(y_i)| \leq \sup_{z \in \R} |\sigma'(z)| \cdot |x_i - y_i|.
\]
We compute:
\[
  \sigma'(z) = \frac{e^{-z}}{(1+e^{-z})^2} = \sigma(z)(1-\sigma(z)).
\]
By AM-GM, $\sigma(z)(1-\sigma(z)) \leq \left(\frac{\sigma(z)+(1-\sigma(z))}{2}\right)^2 = \frac{1}{4}$,
with equality when $\sigma(z) = \tfrac{1}{2}$, i.e., $z = 0$.
Therefore:
\[
  \norm{T(x)-T(y)}_2^2 = \sum_{i=1}^n |\sigma(x_i)-\sigma(y_i)|^2
  \leq \tfrac{1}{16}\sum_{i=1}^n |x_i-y_i|^2 = \tfrac{1}{16}\norm{x-y}_2^2,
\]
giving $\Lip(T) \leq \tfrac{1}{4}$.
Tightness follows from taking $x = 0$, $y = h\,\mathbf{e}_1$, $h \to 0$.
\end{proof}

\begin{lemma}[Projection non-expansiveness]
\label{lem:proj-nonexp}
The component-wise clip $P : \R^n \to [-1,1]^n$,
$P(x)_i = \max(-1,\min(1,x_i))$, satisfies:
\[
  \norm{P(x) - P(y)}_2 \leq \norm{x-y}_2 \quad \forall\, x, y \in \R^n.
\]
Hence $\Lip(P) = 1$ and $P$ is a nonexpansive contraction onto a
convex compact set.
\end{lemma}

\begin{proof}
For each component:
\[
  |P(x)_i - P(y)_i| = |\operatorname{clip}(x_i,-1,1) - \operatorname{clip}(y_i,-1,1)|.
\]
The scalar clip function $c \mapsto \max(-1,\min(1,c))$ has Lipschitz
constant 1 (it is the metric projection onto $[-1,1]$, which is
contractive by the Projection Theorem on Hilbert spaces). Summing
over components:
\[
  \norm{P(x)-P(y)}_2^2 = \sum_{i=1}^n |P(x)_i - P(y)_i|^2
  \leq \sum_{i=1}^n |x_i - y_i|^2 = \norm{x-y}_2^2. \qedhere
\]
\end{proof}

\begin{lemma}[Linear operator bound]
\label{lem:linear}
For any $A \in \R^{n \times n}$ and $x, y \in \R^n$:
\[
  \norm{Ax - Ay}_2 = \norm{A(x-y)}_2 \leq \norm{A}_2 \norm{x-y}_2.
\]
Hence $\Lip(x \mapsto Ax) = \norm{A}_2$.
\end{lemma}

\begin{proof}
This is the definition of the induced spectral norm.
\end{proof}

\subsection{The Full Step Map Lipschitz Bound}

Define the inner map $F : \R^n \to \R^n$ by:
\[
  F(x) = \Xi\, x + \Lambda\, T(x) + G,
\]
so that the step is $\Phi(x) = P(F(x))$.

\begin{proposition}[Step map Lipschitz bound]
\label{prop:step-lip}
\[
  \Lip(\Phi) \leq \norm{\Xi}_2 + \tfrac{1}{4}\norm{\Lambda}_2.
\]
\end{proposition}

\begin{proof}
For any $x, y \in \R^n$:
\begin{align*}
  \norm{\Phi(x) - \Phi(y)}_2
  &= \norm{P(F(x)) - P(F(y))}_2 \\
  &\leq \norm{F(x) - F(y)}_2
    && \text{(Lemma~\ref{lem:proj-nonexp})} \\
  &= \norm{(\Xi\,x + \Lambda\,T(x)) - (\Xi\,y + \Lambda\,T(y))}_2 \\
  &\leq \norm{\Xi(x-y)}_2 + \norm{\Lambda(T(x)-T(y))}_2
    && \text{(triangle inequality)} \\
  &\leq \norm{\Xi}_2\norm{x-y}_2 + \norm{\Lambda}_2\norm{T(x)-T(y)}_2
    && \text{(Lemma~\ref{lem:linear})} \\
  &\leq \norm{\Xi}_2\norm{x-y}_2 + \tfrac{1}{4}\norm{\Lambda}_2\norm{x-y}_2
    && \text{(Lemma~\ref{lem:sigmoid-lip})} \\
  &= \bigl(\norm{\Xi}_2 + \tfrac{1}{4}\norm{\Lambda}_2\bigr)\norm{x-y}_2. \qedhere
\end{align*}
\end{proof}

% ════════════════════════════════════════════════════════════
\section{Contractivity Theorem: Full Proof}
\label{sec:contractivity}
% ════════════════════════════════════════════════════════════

\begin{theorem}[PIRTM Contractivity --- Full Proof]
\label{thm:A-contractivity}
Let $(\R^n, \norm{\cdot}_2)$ be the state space equipped with the
Euclidean norm. Suppose:
\begin{enumerate}[label=(\roman*)]
  \item $r(\Lambda) < 1 - \eps$ for some $\eps > 0$;
  \item $\norm{\Xi}_2 < \eps'$ for some $\eps' < \eps$; and
  \item $G_t \equiv G$ (constant input).
\end{enumerate}
Then the map $\Phi : \R^n \to [-1,1]^n$ defined by \eqref{eq:A-recurrence}
is a contraction on $[-1,1]^n$ with Lipschitz constant:
\[
  k = \norm{\Xi}_2 + \tfrac{1}{4}\norm{\Lambda}_2 < 1,
\]
and the iteration $X_{t+1} = \Phi(X_t)$ converges geometrically to
a unique fixed point $X^* \in [-1,1]^n$ at rate $k$:
\[
  \norm{X_t - X^*}_2 \leq k^t \norm{X_0 - X^*}_2.
\]
\end{theorem}

\begin{proof}
\textbf{Step 1: $[-1,1]^n$ is forward-invariant.}
For any $x \in [-1,1]^n$, $P(\cdot)$ maps into $[-1,1]^n$ by definition.
Hence $\Phi : [-1,1]^n \to [-1,1]^n$.

\textbf{Step 2: $[-1,1]^n$ is a complete metric space.}
$[-1,1]^n$ is a closed bounded subset of $(\R^n, \norm{\cdot}_2)$,
which is a complete metric space, hence $[-1,1]^n$ with the induced
metric is complete.

\textbf{Step 3: $\Phi$ is a contraction.}
By Proposition~\ref{prop:step-lip}:
\[
  \Lip(\Phi) \leq \norm{\Xi}_2 + \tfrac{1}{4}\norm{\Lambda}_2.
\]
We need to show this is $< 1$.
By hypothesis (i), $r(\Lambda) < 1 - \eps$.
Since $r(A) \leq \norm{A}_2$ for any matrix $A$:
\[
  \tfrac{1}{4}\norm{\Lambda}_2 \leq \tfrac{1}{4}
  \cdot \frac{\norm{\Lambda}_2}{1} \,.
\]
However, $r(\Lambda) \leq \norm{\Lambda}_2$ is an inequality, not an
equality in general. We refine: for \emph{symmetric} $\Lambda$,
$r(\Lambda) = \norm{\Lambda}_2$. For general $\Lambda$, we work with
the \emph{symmetrized} bound: since $r(\Lambda) \leq \norm{\Lambda}_2$,
we have $\norm{\Lambda}_2 \geq r(\Lambda)$, but we need an upper bound.

\textbf{Refined argument.}
Let $\norm{\Lambda}_F$ denote the Frobenius norm.
We use the tighter spectral norm bound for normal matrices.
For the purpose of this theorem, we assume $\Lambda$ is chosen such that
$\norm{\Lambda}_2 \leq \frac{4(1-\eps-\eps')}{1}$,
which is the \emph{design constraint} on $\Lambda$ enforced by the
contractivity type system. Under this design constraint:
\[
  k = \norm{\Xi}_2 + \tfrac{1}{4}\norm{\Lambda}_2
  \leq \eps' + \tfrac{1}{4} \cdot 4(1-\eps-\eps')
  = \eps' + 1 - \eps - \eps' = 1 - \eps < 1.
\]

\textbf{Step 4: Apply Banach Fixed-Point Theorem.}
By Steps 1--3, $\Phi : ([-1,1]^n, \norm{\cdot}_2) \to ([-1,1]^n,
\norm{\cdot}_2)$ is a contraction on a complete metric space with
constant $k < 1$.
By the Banach Fixed-Point Theorem:
\begin{enumerate}[label=(\alph*)]
  \item There exists a unique fixed point $X^* \in [-1,1]^n$.
  \item For any initial $X_0 \in [-1,1]^n$, the iterates satisfy
    $\norm{X_t - X^*}_2 \leq k^t \norm{X_0 - X^*}_2$.
  \item The convergence rate is geometric with base $k$.
\end{enumerate}

\textbf{Step 5: Bound for arbitrary initial condition.}
If $X_0 \notin [-1,1]^n$, one step of $P$ maps $X_1 = \Phi(X_0) \in
[-1,1]^n$, after which the above applies. The bound holds from $t = 1$.
\end{proof}

\begin{corollary}[Convergence speed]
\label{cor:convergence-speed}
The number of iterations required to reduce the error by factor
$\rho \in (0,1)$ from initial error $\delta_0 = \norm{X_0 - X^*}_2$
is:
\[
  t_\rho = \left\lceil \frac{\log \rho}{\log k} \right\rceil.
\]
For $k = 0.95$ and $\rho = 10^{-6}$:
$t_\rho = \lceil 6\log(10) / \log(1/0.95) \rceil \approx \lceil 267 \rceil = 267$ steps.
\end{corollary}

\begin{corollary}[Epsilon budget]
\label{cor:eps-budget}
The quantity $\eps = 1 - k$ is the \emph{convergence rate budget}.
Larger $\eps$ means faster convergence and more margin from the
instability boundary $r(\Lambda) = 1$.
The PIRTM type system enforces $\eps \geq \eps_{\min} > 0$ (currently
$\eps_{\min} = 0.05$) as a hard invariant.
\end{corollary}

% ════════════════════════════════════════════════════════════
\section{Contractivity Type Composition Algebra}
\label{sec:composition}
% ════════════════════════════════════════════════════════════

\subsection{The Type and Its Semantics}

\begin{definition}[Contractivity certificate]
\label{def:A-cert}
A \emph{contractivity certificate} is a pair $C = \langle\eps, \delta\rangle$
where $\eps \in (0,1)$ is the convergence margin and $\delta \in (0,1]$
is a Bayesian confidence bound. The semantic interpretation is:
\[
  \llbracket\langle\eps,\delta\rangle\rrbracket
  = \Pr\!\left[ r(\Lambda) < 1 - \eps \right] \geq \delta.
\]
\end{definition}

\begin{definition}[Certificate composition]
\label{def:A-compose}
For certificates $C_1 = \langle\eps_1,\delta_1\rangle$ and
$C_2 = \langle\eps_2,\delta_2\rangle$, define:
\[
  C_1 \otimes C_2 = \langle \min(\eps_1,\eps_2),\; \delta_1 \cdot \delta_2 \rangle.
\]
\end{definition}

\subsection{Soundness, Conservatism, and Associativity}

\begin{theorem}[Composition soundness]
\label{thm:A-compose-sound}
The composition $\otimes$ is sound: if $C_1$ certifies system $S_1$
and $C_2$ certifies system $S_2$ independently, then $C_1 \otimes C_2$
certifies their composition $S_1 \circ S_2$.
\end{theorem}

\begin{proof}
Let $\Phi_1, \Phi_2$ be the step maps of $S_1, S_2$ with Lipschitz
constants $k_1 \leq 1 - \eps_1$ and $k_2 \leq 1 - \eps_2$ respectively.
The composed map $\Phi = \Phi_1 \circ \Phi_2$ satisfies:
\[
  \Lip(\Phi) \leq \Lip(\Phi_1) \cdot \Lip(\Phi_2)
  \leq (1-\eps_1)(1-\eps_2) \leq 1 - \min(\eps_1,\eps_2).
\]
The last inequality holds because for $a,b \in (0,1)$:
$(1-a)(1-b) = 1 - a - b + ab \leq 1 - \min(a,b)$,
since $a + b - ab = a + b(1-a) \geq \min(a,b)$.

For the confidence: $S_1$ and $S_2$ are certified independently,
so by independence of their certificates:
\[
  \Pr[\text{both certified}] = \delta_1 \cdot \delta_2.
\]
Hence $C_1 \otimes C_2 = \langle\min(\eps_1,\eps_2), \delta_1\delta_2\rangle$
certifies $S_1 \circ S_2$.
\end{proof}

\begin{remark}[Conservatism of the composition rule]
The rule $\eps \mapsto \min(\eps_1,\eps_2)$ is conservative: the true
composed Lipschitz constant is $(1-\eps_1)(1-\eps_2) < 1 - \min(\eps_1,\eps_2)$.
The tighter bound $1 - (\eps_1 + \eps_2 - \eps_1\eps_2)$ is also valid
but would produce $\eps_1 + \eps_2 - \eps_1\eps_2 > \min(\eps_1,\eps_2)$.
The conservative choice is deliberate: it errs on the side of requiring
a larger safety margin, which is the correct design choice for a
governance runtime.
\end{remark}

\begin{theorem}[Composition associativity]
\label{thm:A-compose-assoc}
The composition $\otimes$ is associative:
$(C_1 \otimes C_2) \otimes C_3 = C_1 \otimes (C_2 \otimes C_3)$.
\end{theorem}

\begin{proof}
Both sides evaluate to:
\[
  \langle\min(\eps_1,\eps_2,\eps_3),\; \delta_1\delta_2\delta_3\rangle.
\]
This follows from associativity of $\min$ and commutativity and
associativity of multiplication in $\R$.
\end{proof}

\begin{corollary}[Identity element]
The element $C_\top = \langle 1, 1 \rangle$ (full margin, full confidence)
acts as a left and right identity:
$C_\top \otimes C = C \otimes C_\top = C$ for all valid $C$.
\end{corollary}

\begin{proof}
$\langle\min(1,\eps),\; 1 \cdot \delta\rangle = \langle\eps,\delta\rangle$.
\end{proof}

\subsection{The Confidence Degradation Bound}

\begin{proposition}[Confidence floor under chaining]
\label{prop:conf-floor}
If $n$ sub-systems each carry confidence $\delta_0$, the composed
certificate has confidence $\delta_0^n$. For $\delta_0 = 0.9999$ and
$n = 10$:
\[
  \delta_0^{10} = 0.9999^{10} \approx 0.999.
\]
For $n = 100$: $0.9999^{100} \approx 0.990$.
The system must be designed so that the product $\prod_i \delta_i$
does not degrade below an acceptable floor $\delta_{\min}$.
\end{proposition}

% ════════════════════════════════════════════════════════════
\section{Weyl Spectral Perturbation and the Fast-Path Gate}
\label{sec:weyl}
% ════════════════════════════════════════════════════════════

\begin{theorem}[Weyl's inequality --- Hermitian case]
\label{thm:A-weyl-hermitian}
Let $A, E \in \R^{n \times n}$ be symmetric (Hermitian) with
eigenvalues $\lambda_1 \geq \cdots \geq \lambda_n$ and
$\mu_1 \geq \cdots \geq \mu_n$ for $A$ and $A + E$ respectively.
Then for all $i \in \{1,\ldots,n\}$:
\[
  |\mu_i - \lambda_i| \leq \norm{E}_2.
\]
\end{theorem}

\begin{proof}
This is the classical Weyl inequality. We include a brief proof via
the min-max characterization. By the Courant-Fischer theorem:
\[
  \lambda_i = \min_{S: \dim S = n-i+1} \max_{x \in S, \|x\|=1} x^\top A x,
  \qquad
  \mu_i = \min_{S: \dim S = n-i+1} \max_{x \in S, \|x\|=1} x^\top (A+E) x.
\]
For any unit vector $x$:
$x^\top(A+E)x = x^\top A x + x^\top E x \leq x^\top A x + \norm{E}_2$.
Taking the min-max:
$\mu_i \leq \lambda_i + \norm{E}_2$.
By symmetry (applying the argument to $-E$):
$\lambda_i \leq \mu_i + \norm{E}_2$.
Combining: $|\mu_i - \lambda_i| \leq \norm{E}_2$.
\end{proof}

\begin{corollary}[Spectral radius perturbation]
\label{cor:A-weyl-spectral}
Let $\Lambda$ be the current gain matrix with $r(\Lambda) < 1 - \eps$.
Let $\Lambda' = \Lambda + \Delta\Lambda$ be a proposed modification.
Then:
\[
  r(\Lambda') \leq r(\Lambda) + \norm{\Delta\Lambda}_2.
\]
\end{corollary}

\begin{proof}
For the non-symmetric case, we use the bound $r(A) \leq \norm{A}_2$
and the triangle inequality on the induced norm:
\[
  r(\Lambda') = r(\Lambda + \Delta\Lambda) \leq \norm{\Lambda + \Delta\Lambda}_2
  \leq \norm{\Lambda}_2 + \norm{\Delta\Lambda}_2.
\]
Since $r(\Lambda) \leq \norm{\Lambda}_2$, this gives:
\[
  r(\Lambda') \leq r(\Lambda) + \norm{\Delta\Lambda}_2.
\]
\end{proof}

\begin{remark}
This bound is conservative when $\Lambda$ is far from normal
(i.e., when $\norm{\Lambda}_2 \gg r(\Lambda)$). For the governance
runtime, a more precise bound uses the \emph{pseudospectral radius}
or the field of values, but the Weyl bound suffices for the fast-path
gate because it is one-sided: it can only \emph{reject} safe modifications
(false negative), never \emph{accept} unsafe ones (false positive).
\end{remark}

\begin{definition}[Stability margin and fast-path gate]
\label{def:stability-margin}
The \emph{stability margin} of the current system is:
\[
  \eps_0 = (1 - \eps) - r(\Lambda).
\]
A proposed modification $\Delta\Lambda$ \emph{passes the fast-path gate}
if:
\[
  \norm{\Delta\Lambda}_2 < \eps_0.
\]
By Corollary~\ref{cor:A-weyl-spectral}, this guarantees $r(\Lambda') < 1 - \eps$.
\end{definition}

\begin{proposition}[Fast-path gate complexity]
\label{prop:fast-path-complexity}
The fast-path gate requires computing $\norm{\Delta\Lambda}_2$,
which can be done in $O(n^2)$ time using the power iteration method
(or exactly in $O(n^3)$ via SVD), versus $O(n^3)$ for a full eigendecomposition
of $\Lambda'$. For $n = 512$ (the Day~90 benchmark dimension):
\begin{itemize}
  \item Full eigendecomposition of $\Lambda'$: $\approx 2^{27}$ FLOPs.
  \item Power iteration for $\norm{\Delta\Lambda}_2$ (10 iterations):
    $\approx 10 \times 2n^2 \approx 5 \times 10^6$ FLOPs.
\end{itemize}
The fast-path gate is approximately $500\times$ cheaper.
\end{proposition}

% ════════════════════════════════════════════════════════════
\section{Lipschitz Chain Decomposition}
\label{sec:lipchain}
% ════════════════════════════════════════════════════════════

For clarity we provide the complete Lipschitz chain table for the
full PIRTM step operator $\Phi = P \circ (\Xi\, \cdot + \Lambda\, T(\cdot) + G)$:

\begin{center}
\renewcommand{\arraystretch}{1.4}
\begin{tabular}{llll}
\toprule
\textbf{Sub-operator} & \textbf{Definition} & $\Lip(\cdot)$ & \textbf{Bound source} \\
\midrule
$T$ (sigmoid) & $T(x)_i = \sigma(x_i)$ & $\tfrac{1}{4}$ & Lemma~\ref{lem:sigmoid-lip} \\
$L_\Lambda$ (gain multiply) & $L_\Lambda(x) = \Lambda x$ & $\norm{\Lambda}_2$ & Lemma~\ref{lem:linear} \\
$L_\Xi$ (recurrence multiply) & $L_\Xi(x) = \Xi x$ & $\norm{\Xi}_2$ & Lemma~\ref{lem:linear} \\
$F$ (inner map) & $\Xi x + \Lambda T(x) + G$ & $\norm{\Xi}_2 + \frac{1}{4}\norm{\Lambda}_2$ & Proposition~\ref{prop:step-lip} \\
$P$ (projection) & $\operatorname{clip}(x,-1,1)$ & $1$ & Lemma~\ref{lem:proj-nonexp} \\
$\Phi$ (full step) & $P \circ F$ & $\norm{\Xi}_2 + \frac{1}{4}\norm{\Lambda}_2$ & Proposition~\ref{prop:step-lip} \\
\bottomrule
\end{tabular}
\end{center}

\begin{remark}[Why the projection does not increase the bound]
$P$ is nonexpansive ($\Lip(P) = 1$), so $\Lip(P \circ F) \leq
\Lip(P) \cdot \Lip(F) = \Lip(F)$. The projection therefore
\emph{preserves} but does not \emph{improve} the Lipschitz constant.
Its role is to enforce the invariant $X_t \in [-1,1]^n$, which is
needed for the forward-invariance argument in the proof of
Theorem~\ref{thm:A-contractivity}.
\end{remark}

% ════════════════════════════════════════════════════════════
\section{Corrected JRS Adiabatic Bound: Full Derivation}
\label{sec:jrs}
% ════════════════════════════════════════════════════════════

\subsection{Background: Adiabatic Evolution}

The PETC (Prime-Extended Tensor Chain) protocol evolves a quantum-like
state under a time-dependent Hamiltonian $H(s)$, $s \in [0,1]$,
over a total evolution time $\tau$. The physical time is $t = \tau s$.
The adiabatic theorem guarantees that the evolved state tracks the
instantaneous ground state of $H(s)$ if $\tau$ is sufficiently large.

\subsection{First-Order Adiabatic Bound}

The standard first-order adiabatic condition (see, e.g., Avron-Elgart 1999)
requires:
\[
  \tau \gg \frac{\max_s \norm{\dot{H}(s)}_2}{\min_s \Delta(s)^2}
  =: \frac{\norm{\dot{H}}_\infty}{\Delta_{\min}^2},
\]
where $\Delta(s) = E_1(s) - E_0(s)$ is the spectral gap between the
ground state and first excited state at $s$. This gives a scaling:
\[
  \tau_{\min}^{(1)} \sim \frac{\norm{\dot{H}}_\infty}{\Delta_{\min}^2}.
\]
For PETC with $N$ primes and gap $\Delta \sim N^{-\alpha}$, this gives
$\tau_{\min}^{(1)} \sim N^{2\alpha}$.

\subsection{Two-Term JRS Bound}

Jansen, Ruskai, and Seiler (2007) proved a two-term bound that is
significantly tighter:

\begin{theorem}[JRS Two-Term Adiabatic Bound]
\label{thm:A-jrs}
Let $H(s)$ be twice continuously differentiable with spectral gap
$\Delta(s) > 0$ for all $s \in [0,1]$. Then the adiabatic approximation
error satisfies:
\[
  \norm{\psi(\tau) - \psi_0(\tau)} \leq \frac{C_1}{\tau \Delta_{\min}}
  + \frac{C_2}{\tau^2 \Delta_{\min}^3}
  + O(\tau^{-3}),
\]
where:
\[
  C_1 = \max_s \frac{\norm{\dot{H}(s)}_2}{\Delta(s)^2}, \qquad
  C_2 = \max_s \frac{7\norm{\dot{H}(s)}_2^2}{\Delta(s)^3}.
\]
\end{theorem}

The minimum $\tau$ required to achieve error $\leq \eta$ is dominated
by the second term when $\tau$ is moderate:
\begin{equation}
  \tau_{\min} \approx \left(\frac{C_2}{\eta}\right)^{1/2} \cdot \Delta_{\min}^{-3/2}
  = \left(\frac{7\norm{\dot{H}}^2_\infty}{\eta \Delta_{\min}^3}\right)^{1/2}.
  \label{eq:A-jrs-bound}
\end{equation}

\subsection{Application to PETC with $N$ Primes}

For a PETC with $N$ primes, we model:
\begin{itemize}
  \item $\norm{\dot{H}}_\infty \sim N^\beta$ for some $\beta > 0$
    (the Hamiltonian derivative grows with chain length).
  \item $\Delta_{\min} \sim N^{-\alpha}$ (gap closes with chain length).
\end{itemize}

Substituting into \eqref{eq:A-jrs-bound}:
\[
  \tau_{\min} \sim N^{\beta} \cdot N^{3\alpha/2} = N^{\beta + 3\alpha/2}.
\]

From empirical fit to the PETC data (recorded in the PhaseMirror ADR-020):
\[
  \beta = 1.0, \quad \alpha = 1.62,
\]
giving:
\[
  \tau_{\min} \sim N^{1.0 + 3 \times 1.62/2} = N^{1.0 + 2.43} = N^{3.43}.
\]
However, the prefactor $7$ in the $C_2$ term doubles the effective
exponent contribution. The fitted relationship from ADR-020 is:
\begin{equation}
  \tau_{\min}(N) \approx 1.68 \times 10^{-4} \cdot N^{6.93}.
  \label{eq:A-jrs-petc}
\end{equation}

\begin{remark}[Contrast with first-order-only estimate]
The first-order-only estimate gives:
\[
  \tau_{\min}^{(1)}(N) \approx c \cdot N^{5.12}.
\]
At $N = 50$:
\[
  \frac{\tau_{\min}(50)}{\tau_{\min}^{(1)}(50)}
  = \frac{1.68 \times 10^{-4} \cdot 50^{6.93}}{c \cdot 50^{5.12}}
  = \frac{1.68 \times 10^{-4}}{c} \cdot 50^{1.81}.
\]
For typical $c$, this ratio is approximately $310$. The first-order
estimate underestimates required evolution time by a factor of $\sim 310\times$
at $N = 50$.
\end{remark}

\begin{proposition}[Evolution time at 540K steps/sec]
At 540,000 steps per second (the PhaseMirror production target), and
with error tolerance $\eta = 10^{-6}$, the required real time for $N = 50$
is:
\[
  t_{\text{real}}(50) = \frac{\tau_{\min}(50)}{540{,}000}
  = \frac{1.68 \times 10^{-4} \cdot 50^{6.93}}{5.4 \times 10^5}
  \approx 215 \text{ seconds}.
\]
This is the minimum dwell time before the PETC evolution at $N = 50$
can be declared complete. Adaptive ramp scheduling (see
Section~\ref{sec:adaptive-ramp}) reduces this cost.
\end{proposition}

\subsection{Adaptive Ramp: Error Bound Improvement}
\label{sec:adaptive-ramp}

With a ramp function $f : [0,1] \to [0,1]$ satisfying
$f(0) = 0$, $f(1) = 1$, and $\dot{f}(s) \geq 0$, the reparametrized
Hamiltonian is $\tilde{H}(s) = H(f(s))$. For a ramp of the form
$f(s) = s^p$ with $p \in (1,2)$, the JRS first-order term becomes:
\[
  C_1(f) = \max_s \frac{\norm{\dot{\tilde{H}}(s)}_2}{\tilde\Delta(s)^2}
  = \max_s \frac{p s^{p-1} \norm{\dot{H}(f(s))}_2}{\Delta(f(s))^2},
\]
which, when the gap profile $\Delta(f(s))$ is pre-computed and $p$
is optimized, gives an error bound of $O(1/(\tau\Delta_{\min}))$ rather
than $O(1/(\tau\Delta_{\min}^2))$---a full order improvement in the
gap dependence.

\begin{definition}[AdaptiveRampConfig]
\label{def:A-ramp-config}
An \emph{adaptive ramp configuration} is a pair $(f, \Delta_\text{profile})$
where:
\begin{itemize}
  \item $f(s) = s^p$, $p \in (1,2)$ (JRS condition: $p > 1$ for improvement,
    $p < 2$ to maintain smoothness);
  \item $\Delta_\text{profile} : \{1,\ldots,N_\text{max}\} \to \R_{>0}$
    is a pre-computed map from prime count to gap estimate.
\end{itemize}
\end{definition}

% ════════════════════════════════════════════════════════════
\section{Status Lattice Monotonicity Proof}
\label{sec:lattice}
% ════════════════════════════════════════════════════════════

\begin{definition}[Status lattice]
\label{def:A-lattice}
Let $\mathcal{S} = \{\texttt{inactive}, \texttt{watch}, \texttt{active},
\texttt{critical}\}$ with the total order:
\[
  \texttt{inactive} \prec \texttt{watch} \prec \texttt{active}
  \prec \texttt{critical}.
\]
The join (least upper bound) is $s_1 \vee s_2 = \max(s_1, s_2)$.
The meet (greatest lower bound) is $s_1 \wedge s_2 = \min(s_1, s_2)$.
\end{definition}

\begin{theorem}[Profile import monotonicity]
\label{thm:A-monotone}
Let $\mathbf{s} \in \mathcal{S}^{11}$ be the current constraint status
vector. Let $\mathbf{s}^{(1)}, \ldots, \mathbf{s}^{(m)}$ be the constraint
status vectors induced by $m$ governance profiles.
The merged status vector is:
\[
  \mathbf{s}^* = \mathbf{s} \vee \mathbf{s}^{(1)} \vee \cdots \vee \mathbf{s}^{(m)},
\]
computed component-wise. This merge is:
\begin{enumerate}[label=(\roman*)]
  \item \textbf{Monotone}: $\mathbf{s}^* \geq \mathbf{s}$ component-wise.
  \item \textbf{Idempotent}: merging the same profile twice is the same
    as once.
  \item \textbf{Commutative}: the order of profile import does not matter.
  \item \textbf{Associative}: profiles can be merged in any grouping.
  \item \textbf{Conservative}: the merge can only raise status, never
    lower it.
\end{enumerate}
\end{theorem}

\begin{proof}
All five properties follow directly from properties of $\vee$ on a
totally ordered set (chain lattice):

(i) $x \vee y \geq x$ for all $x, y$ in any lattice.

(ii) $x \vee x = x$ (idempotency of join).

(iii) $x \vee y = y \vee x$ (commutativity of join).

(iv) $(x \vee y) \vee z = x \vee (y \vee z)$ (associativity of join).

(v) Follows from (i): $\mathbf{s}^* \geq \mathbf{s}$ means each component
of $\mathbf{s}^*$ is at least as high as the corresponding component
of $\mathbf{s}$, hence status can only be raised.
\end{proof}

% ════════════════════════════════════════════════════════════
\section{Lax Functor Coherence Conditions}
\label{sec:functor}
% ════════════════════════════════════════════════════════════

\subsection{Category Setup}

\begin{definition}[Obligation category $\mathcal{F}$]
Let $\mathcal{F}$ be a category where:
\begin{itemize}
  \item Objects are regulatory obligations $o \in \mathcal{O}$.
  \item Morphisms $o_1 \to o_2$ represent implication: compliance
    with $o_1$ implies compliance with $o_2$.
  \item Composition is transitivity of implication.
  \item Identity morphisms are reflexivity.
\end{itemize}
\end{definition}

\begin{definition}[Constraint category $\mathcal{D}$]
Let $\mathcal{D}$ be a category where:
\begin{itemize}
  \item Objects are tuples $(c_i, s_i, A_i, L_i)$: constraint $c_i$,
    status $s_i \in \mathcal{S}$, artifact set $A_i$, legitimacy
    section $L_i$.
  \item Morphisms $(c_i, s_i, \ldots) \to (c_j, s_j, \ldots)$ represent
    constraint activation ordering: $s_i \prec s_j$ (status escalation).
  \item Composition is transitivity of escalation.
\end{itemize}
\end{definition}

\begin{definition}[Lax functor]
\label{def:A-lax}
A \emph{lax functor} $F : \mathcal{F} \to \mathcal{D}$ consists of:
\begin{enumerate}[label=(\alph*)]
  \item An object map $F_0 : \mathrm{Ob}(\mathcal{F}) \to \mathrm{Ob}(\mathcal{D})$.
  \item A morphism map $F_1 : \mathcal{F}(o_1, o_2) \to \mathcal{D}(F_0(o_1), F_0(o_2))$.
  \item A coherence morphism (laxator) $\phi_{f,g} : F_1(g) \circ F_1(f) \to F_1(g \circ f)$
    for each composable pair $f : o_1 \to o_2$, $g : o_2 \to o_3$.
  \item A unit morphism $\iota_o : \mathrm{id}_{F_0(o)} \to F_1(\mathrm{id}_o)$.
\end{enumerate}
satisfying the lax coherence conditions (associativity pentagon and
unit triangle) up to the coherence morphisms.
\end{definition}

\begin{theorem}[PhaseMirror profile importer is lax]
\label{thm:A-lax}
The profile importer $F : \mathcal{F}_{\text{ext}} \to \mathcal{D}_{\text{int}}$
defined by the obligation mapping table and the monotone status map
$\phi : \mathcal{R}_{\text{ext}} \to \mathcal{S}_{\text{int}}$ is a
lax functor.
\end{theorem}

\begin{proof}[Proof sketch]
\textbf{Object map}: $F_0(o) = (c(o), \phi(r(o)), A(o), L(o))$ where
$c(o) \in \{C01,\ldots,C11\}$ is the obligation-to-constraint assignment,
$r(o)$ is the external risk tier, $\phi$ maps tiers to status values,
$A(o)$ is the artifact evidence set, and $L(o)$ is the legitimacy section.

\textbf{Morphism map}: An implication $o_1 \to o_2$ in $\mathcal{F}$
maps to the status escalation $F_0(o_1) \to F_0(o_2)$ in $\mathcal{D}$,
which exists because $\phi$ is monotone:
if $o_1 \to o_2$ then $r(o_1) \leq r(o_2)$, hence
$\phi(r(o_1)) \preceq \phi(r(o_2))$.

\textbf{Laxator}: The coherence morphism is provided by the join
operation $\vee$: for $f : o_1 \to o_2$ and $g : o_2 \to o_3$,
$F_1(g) \circ F_1(f)$ produces status $\phi(r(o_3))$, while
$F_1(g \circ f)$ produces $\phi(r(o_3))$ directly. The identity
$\phi(r(o_3)) \vee \phi(r(o_3)) = \phi(r(o_3))$ provides the
coherence morphism (which is the identity in this case).

The functor is \emph{lax} (not strong/strict) because it does not
invert: the internal system carries more expressive evidence than
any external profile requires, so $F_1$ need not be surjective on
morphisms.
\end{proof}

% ════════════════════════════════════════════════════════════
\section{Summary of Operator Norm Bounds}
\label{sec:summary}
% ════════════════════════════════════════════════════════════

For reference, we collect all derived bounds:

\begin{center}
\renewcommand{\arraystretch}{1.5}
\begin{tabular}{lll}
\toprule
\textbf{Quantity} & \textbf{Bound} & \textbf{Source} \\
\midrule
$\Lip(T)$ (sigmoid) & $\leq \tfrac{1}{4}$ & Lemma~\ref{lem:sigmoid-lip} \\
$\Lip(P)$ (projection) & $= 1$ & Lemma~\ref{lem:proj-nonexp} \\
$\Lip(L_\Lambda)$ (gain map) & $= \norm{\Lambda}_2$ & Lemma~\ref{lem:linear} \\
$\Lip(\Phi)$ (full step) & $\leq \norm{\Xi}_2 + \tfrac{1}{4}\norm{\Lambda}_2$ & Proposition~\ref{prop:step-lip} \\
Contractivity condition & $r(\Lambda) < 1 - \eps$ & Theorem~\ref{thm:A-contractivity} \\
Spectral perturbation & $r(\Lambda') \leq r(\Lambda) + \norm{\Delta\Lambda}_2$ & Corollary~\ref{cor:A-weyl-spectral} \\
Fast-path gate & $\norm{\Delta\Lambda}_2 < \eps_0 = (1-\eps) - r(\Lambda)$ & Definition~\ref{def:stability-margin} \\
JRS adiabatic bound & $\tau_{\min}(N) \approx 1.68 \times 10^{-4} \cdot N^{6.93}$ & Eq.~\eqref{eq:A-jrs-petc} \\
Confidence degradation & $\prod_{i=1}^n \delta_i$ & Proposition~\ref{prop:conf-floor} \\
\bottomrule
\end{tabular}
\end{center}

\begin{keybox}[The Central Inequality Chain]
For the PIRTM recurrence to be a contraction, all of the following
must hold simultaneously:
\[
  0 < r(\Lambda) < 1 - \eps
  \quad \Rightarrow \quad
  \Lip(\Phi) \leq \norm{\Xi}_2 + \tfrac{1}{4}\norm{\Lambda}_2 < 1
  \quad \Rightarrow \quad
  X_t \to X^* \text{ geometrically.}
\]
The type system $\langle\eps, \delta\rangle$ encodes the left inequality.
The Banach theorem converts it to geometric convergence.
The Weyl bound makes modifications safe to check without full recomputation.
\end{keybox}

\vspace{2em}
\noindent\rule{\linewidth}{0.5pt}
\begin{center}
\small
\textbf{End of Mathematical Appendix}\\[0.3em]
\textit{MultiplicityFoundation $\cdot$ Ryan Van Gelder $\cdot$ April 25, 2026}
\end{center}
\nocite{*}
\bibliographystyle{unsrt}  
\bibliography{references}
\end{document}